\documentclass[10pt,a4paper]{amsart}
\usepackage[margin=19mm]{geometry}
\usepackage{amsmath,amssymb,mathtools}
\usepackage[scr=rsfs,cal=euler]{mathalfa}
\usepackage{xfrac}
\usepackage[unicode,hidelinks,bookmarks=false]{hyperref}
\newcommand{\ie}{i.e.\ }
\newcommand{\cf}{cf.\ }
\newcommand{\resp}{respectively\ }
\newcommand{\Subsection}{\subsection}
\providecommand{\nobreakdash}{\nobreak\hbox{-}}
\setlength{\emergencystretch}{2em}
\multlinegap0pt
\usepackage{bm}
\usepackage{bbm}
\usepackage{dsfont}
\usepackage{xspace}
\usepackage{cancel}
\usepackage{mathtools}
\usepackage{cleveref}
\usepackage[shortlabels]{enumitem}
\usepackage{amsthm}
\makeatletter
\@ifundefined{theorem}{\newtheorem{theorem}{Theorem}}{}
\@ifundefined{lemma}{\newtheorem{lemma}[theorem]{Lemma}}{}
\@ifundefined{property}{\newtheorem{property}{Property}}{}
\makeatother
\newcommand{\Inv}{\operatorname{Inv}}
\newcommand{\car}{\mathbin{\Box}}
\newcommand{\dist}{\operatorname{dist}}
\newcommand{\supp}{\operatorname{supp}}
\renewcommand{\emptyset}{\varnothing}
\renewcommand{\geq}{\geqslant}
\renewcommand{\leq}{\leqslant}
\renewcommand{\theproperty}{P\arabic{property}}
\title[Majority bootstrap percolation on the permutahedron and othe]{Majority bootstrap percolation on the permutahedron and other high-dimensional graphs}
\author{Maur\'icio Collares, Joshua Erde, Anna Geisler, Mihyun Kang}
\date{Source: arXiv:2507.06597v1 (CC BY 4.0)}
\begin{document}
\maketitle

\noindent\emph{Source and licence.} Majority bootstrap percolation on the permutahedron and other high-dimensional graphs. Version arXiv:2507.06597v1 (CC BY 4.0), distributed under the Creative Commons Attribution 4.0 International licence, as recorded in the arXiv licence metadata for this exact version and shown on the abstract page https://arxiv.org/abs/2507.06597v1. Full rights notice: licenses/arxiv-2507.06597-CC-BY-4.0.txt.\par\smallskip

\noindent\textit{Editorial scope.} Counted unit: the Lemma whose source label is \texttt{l:permutahedronH} in the article (source0.tex lines 1173--1177, source label \texttt{l:permutahedronH}), delivered together with the article's own complete original proof of it (source0.tex lines 1179--1278, one \texttt{proof} environment of 100 lines). No mathematics is written, repaired or re-derived: statement and proof are the article's own text line for line. Required context retained uncounted: c:regular (lines 400--402); c:backwards (lines 404--409); c:projection (lines 425--432); c:partitioncond (lines 434--439); c:localconnectionstrong (lines 453--462); i:backwardstrong (lines 455--461); i:cherries (lines 455--461); i:cherrynew (lines 455--461); i:smallnew (lines 455--461); i:sparsenew (lines 455--461); c:bound-larger (lines 512--514). Every definition, display and label of those lines is retained unchanged. Omitted from this package and not redistributed: 1--399, 403--403, 410--424, 433--433, 440--452, 462--511, 515--1172, 1178--1178, 1279--1532. Nothing inside a retained excerpt is omitted or altered: every retained body is delivered exactly as the article's lines are written, so no omission receipt of any kind accompanies this package. Citations: the retained excerpts carry no citation command at all, so no bibliography entry of the article is transcribed and no reference is redistributed. Reference adaptation: none was needed and none was made; every reference inside the delivered unit resolves inside it, so no unresolved reference or citation is produced. Printed slips kept literal: no printed word, symbol, hypothesis or number of the counted statement or of its proof was altered. Layout and class: the article's own class is \texttt{amsart} with packages including fullpage, hyperref, cleveref, amssymb, mathtools, ifthen, bm, comment, soul, csquotes, tikz, tikz-3dplot, lmodern, enumitem; the wrapper is the lane's pinned \texttt{amsart} editorial wrapper at 10pt on A4, which restates the theorem-like environments the retained text needs and adds the article's own macro definitions that the retained text needs (\texttt{\textbackslash Inv}, \texttt{\textbackslash car}, \texttt{\textbackslash dist}, \texttt{\textbackslash emptyset}, \texttt{\textbackslash geq}, \texttt{\textbackslash leq}, \texttt{\textbackslash supp}, \texttt{\textbackslash theproperty}), with extra wrapper packages (bm, bbm, dsfont, xspace, cancel, mathtools, cleveref, enumitem). The delivered numbering is the unit's own. Assets: no retained span contains a figure environment, a graphics-inclusion command, a PDF-page inclusion, a table or a TikZ picture; no image file is copied into this package, the package declares no asset dependency and no image file is redistributed with it. Licence: version arXiv:2507.06597v1 of this article is distributed under the Creative Commons Attribution 4.0 International licence, as recorded in the arXiv licence metadata for this exact version and shown on the abstract page https://arxiv.org/abs/2507.06597v1. A full rights notice is delivered at licenses/arxiv-2507.06597-CC-BY-4.0.txt. Characters: every retained excerpt is pure ASCII, so the delivered engine, XeLaTeX with its default Latin Modern text font, reports no missing character.
\begin{property}[Locally almost regular]\makeatletter\edef\@currentlabel{(\theproperty)}\makeatother\label{c:regular}
  For every $x \in V(G)$, $\ell \in \mathbb{N}$ and $y \in S(x, \ell)$, $$|d(x)-d(y)| \leq K \ell.$$
\end{property}

\begin{property}[Bounded backwards expansion]\makeatletter\edef\@currentlabel{(\theproperty)}\makeatother\label{c:backwards}
  For every $x \in V(G)$, $\ell \in \mathbb{N}$ and $y \in S(x, \ell)$,
  \[
    |N(y) \cap B(x, \ell)| \leq K\ell.
  \]
\end{property}

\begin{property}[Projection]\makeatletter\edef\@currentlabel{(\theproperty)}\makeatother\label{c:projection}
  For every $x \in V(G)$, $\ell \in \mathbb{N}$ and $y \in S(x, \ell)$, there is a subgraph $G(y)$ of $G$ such that the following hold:
  \begin{enumerate}[$(i)$]
  \item $y \in V(G(y))$;
  \item $G(y) \in \mathcal{F}(K)$;
  \item $V(G(y)) \cap B(x, \ell-1) = \emptyset$; \item $ |d_{G(y)}(w) - d_{G}(w)| \leq K \ell$ for all $w \in V(G(y))$.
  \end{enumerate}
\end{property}

\begin{property}[Separation]\makeatletter\edef\@currentlabel{(\theproperty)}\makeatother\label{c:partitioncond}
  For every $x \in V(G)$, $\ell \in \mathbb{N}$ and $y \in S_0(x, \ell)$, we have
  \[
    |B(y, 2\ell-1) \cap S_0(x, \ell)| \leq \ell K^{\ell-1}d(x)^{\ell-1}.
  \]
\end{property}

\begin{property}[Stronger typical local structure]\makeatletter\edef\@currentlabel{(\theproperty)}\makeatother\label{c:localconnectionstrong}
  For every $x\in V(G)$ there is a set $D \subseteq V(G) \setminus \{x\}$ of \emph{non-typical} vertices such that for every $\ell \in \mathbb{N}$ the following hold:
  \begin{enumerate}[$(i)$]
  \item\label{i:smallnew} $|D\cap S(x,\ell)| \leq  K^{\ell-1} (d(x))^{\ell-1}$;
  \item\label{i:sparsenew} $|D\cap N(y)| \leq K\ell$ for every vertex $y \in S(x, \ell) \setminus D$;
  \item\label{i:cherrynew} every two vertices in $S(x, \ell) \setminus D$ have at most one common neighbour in $S(x, \ell+1) \setminus D$;
  \item\label{i:backwardstrong} every vertex $v \in S(x, \ell) \setminus D$ has at most $\ell$ neighbours in $S(x, \ell-1)$;
  \item\label{i:cherries} for every cherry $uvw$ with $u, w \in S(x, \ell) \setminus D$ and $v \in S(x, \ell+1) \setminus D$, it holds that $(S(x, \ell-1) \setminus D) \cap N(u) \cap N(w) \neq \varnothing$.
  \end{enumerate}
\end{property}

\begin{property}[Larger bounded order]\makeatletter\edef\@currentlabel{(\theproperty)}\makeatother\label{c:bound-larger}
  $G$ has at most $\exp\left(\frac{\delta(G)^{3/2} \log \log \delta(G)}{\log^2 \delta(G)}\right)$ vertices.
\end{property}

\begin{lemma}\label{l:permutahedronH}
  \[
    \mathcal{P}(n) \subseteq \mathcal{H}(4).
  \]
\end{lemma}

\begin{proof}
  We induct on $n$.
  Let $G \in \mathcal{P}(n)$, where by the discussion above we may without loss of generality assume that $G$ is the Cayley graph of a subgroup $H$ of some permutation group $S_{N+1}$ generated by some subset $I$ of the adjacent transpositions with $|I|=n$ which contains the identity.

  \begin{proof}[\Cref{c:regular}]\let\qed\relax
    $G$ is $n$-regular.
  \end{proof}

  \begin{proof}[\Cref{c:backwards}]\let\qed\relax
    This property is preserved under taking isometric subgraphs and is increasing in $K$. Since \Cref{c:backwards} holds in the hypercube $Q^n$ with $K=2$, it is also true in $G$ with $K=4$.
  \end{proof}

  \begin{proof}[\Cref{c:localconnectionstrong}]\let\qed\relax
    For a permutation $\sigma$ of $[n+1]$ let the support of $\sigma$ be the set $\supp(\sigma)=\{i \in [n+1] \,:\, \sigma(i)\neq i\}$.
    Given $x \in V(G)$ and $y \in S(x, \ell)$ for some $\ell \geq 1$, define
    \[
      T(x,y) = (\tau_1,\dots, \tau_{\ell})
    \]
    to be an arbitrary sequence of exactly $\ell$ adjacent transpositions in $I$ such that $y = x\tau_1 \dots \tau_\ell$.
    Set
    \[
      S_0(x, \ell)= \{y \in S(x,\ell) \,:\, \text{all transpositions in $T(x,y)$ have pairwise disjoint supports}\}.
    \]
    By definition $S_0(x,\ell) \subseteq S(x, \ell)$. Note that this definition of $S_0(x,\ell)$ implicitly defines the set
    \[
      D \coloneqq \bigcup_{\ell} S(x,\ell) \setminus S_0(x, \ell).
    \]
    Whenever $y \in S_0(x, \ell)$, the transpositions in the sequence $T(x, y)$ can be applied in any order.
    In this case, we will (in a slight abuse of notation) consider $T(x,y)$ as a subset of $I$, which satisfies $|T(x,y)| = \dist(x,y)=\ell$.


    If $w \in S(x,\ell) \cap D$, then $T(x, w) = (\tau_1, \ldots, \tau_\ell)$ is such that there exist $1 \leq i < j \leq \ell$ with $\supp(\tau_i)\cap\supp(\tau_j)\neq\emptyset$.
    Therefore,
    \[
      |S(x,\ell) \cap D| \leq 3 \binom{\ell}{2} n^{\ell-1},
    \]
    and so \ref{i:smallnew} holds.

    If $w \in S_0(x, \ell)$, then a neighbour of $w$ in $S(x, \ell+1) \cap D$ must differ from $w$ in a transposition that has non-disjoint support with one of the transpositions in $T(x, w)$, and hence there are at most $3 \ell$ of them.
    Note that any $w \in S_0(x, \ell)$ has no neighbours in $S(x, \ell-1) \cap D$, and no neighbours at all in $S(x,\ell)$, so \ref{i:sparsenew} holds.

    Furthermore, since properties \ref{i:cherrynew} and \ref{i:backwardstrong} hold in the hypercube $Q^n$, even when $K=2$ and $D=\emptyset$, and are preserved under taking isometric subgraphs and increasing in $K$, they are also satisfied in $G$ for our choice of $D$ and $K$.

    Finally, let $uvw$ be a cherry with $u, w \in S_0(x, \ell)$ and $v \in S_0(x, \ell+1)$.
    It is easy to verify that $|T(x, u) \cup T(x, w)|=\ell+1$ and $v= u \tau_1 = w \tau_2$ for some adjacent transpositions $\tau_1, \tau_2$ with disjoint support.
    Then $v\tau_1\tau_2$ is a common neighbour of $u$ and $w$ in $S_0(x, \ell-1)$, and so \ref{i:cherries} holds.
\end{proof}

\begin{proof}[\Cref{c:projection}]\let\qed\relax
  Without loss of generality, suppose that $x \in V(G)$ is the identity and let $\ell \in \mathbb{N}$ and $y \in S(x,\ell)$.
  Consider the set
  \[
    \Inv'(y) \coloneqq \{\{i, j\} \subseteq [n+1] \,:\, (i-j)\cdot(y(i)-y(j)) < 0\},
  \]
  that is, the set of pairs of \emph{positions} corresponding to inversions of $y$.
  Let $I(y) \subseteq I$ be the set of generators whose support intersects a pair in $\Inv'(y)$.
  Note that $|I(y)| \leq 4 \ell$.

  Let us consider the subgroup $H(y)$ of $S_{N+1}$ generated by all transpositions except those in $I(y)$ and its Cayley graph $\Gamma(y)$.
  The cosets of $H(y)$ in $\Gamma$ form a decomposition of the vertex set, giving rise to a natural packing of copies of $\Gamma(y)$ inside $G$.

  Let $G(y)$ be the copy of $\Gamma(y)$ which contains $y$.
  Clearly $y \in G(y) \subseteq G$.
  By construction, every vertex $z \in G(y)$ satisfies $\Inv'(y) \subseteq \Inv'(z)$ and hence $G(y) \cap B(x, \ell-1)=\emptyset$ and $G(y) \in \mathcal{P}(n')$ for some $n' \geq n-4\ell$.
  In particular, $G(y)$ is $n'$-regular, and for any $w \in V(G(y))$,
  \[
    |d_{G(y)}(w) - d_G(w)|  = n-n' \leq 4\ell.
  \]

  Finally, by the induction hypothesis, $G(y) \in \mathcal{H}(4)$.
\end{proof}

\begin{proof}[\Cref{c:partitioncond}]\let\qed\relax
  Given $x \in V(G)$, $\ell \in \mathbb{N}$ and $w \in S_0(x,\ell)$, let
  \[
    Z = \{ z \in S_0(x,\ell) \,:\, \dist(z,w) \leq 2\ell-1\}.
  \]
  By considering the isometric embedding of $G$ into the hypercube described above, it can be seen that for any $w,z \in S_0(x,\ell)$
  \[
    \dist(w, z) = |T(x, z) \triangle T(x,w)|=2|T(x,z) \setminus T(x,y)|.
  \]
  Hence, for every $z \in Z$,
  \[
    |T(x,z) \setminus T(x,w)| = \frac{1}{2} \dist(w, z) < \ell.
  \]
  It follows that
  \[
    Z \subseteq Z' = \{ z \in S_0(x, \ell) \,:\, |T(x,z) \setminus T(x,w)| \leq \ell-1\},
  \]
  and it is clear that
  \[
    |Z| \leq |Z'| \leq \sum_{i=0}^{\ell-1} \binom{n}{i} \leq 2n^{\ell-1}.
  \]
  \end{proof}

  \begin{proof}[\Cref{c:bound-larger}]\let\qed\relax
    Finally, if $G= \car_{i=1}^m P_{n_i}$ with $\sum_{i=1}^m n_i =n$, then $|V(G)| = \prod_{i=1}^m (n_i)! \leq n! \leq \exp( n \log n)$.
  \end{proof}
  This concludes the proof.
\end{proof}

\end{document}
